\section{Future Work}
While the software implementation for ECG compressed sensing proposed in this paper yielded adequate results, there is still considerable room for improvement.
Most notably, the results achieved a lower PRD than what has been shown to be possible in the current literature \citep{daponte2021}.
The most significant factor in this was the selection of algorithm for reconstruction.
Due to time constraints, this implementation was limited to convex problems which could be easily solved with the cvxpy library.
Future revisions will focus on adapting more advanced algorithms, such as M-FOCUSS or M-SBL, either by adapting the project to be compatible with existing libraries (cr-sparse) or by writing a custom implementation.
\newline
\par
Furthermore, recent developments in compression of ECG signals have begun to transition away from compressed sensing in favor of deep-learning.
With compressed sensing, the results of recovery depends highly on the assumption that the selected wavelet basis drives sparsity for the measured signal.
As shown by the high variance in PRD when analyzing different sources, this assumption becomes brittle when generalized to a wide variety of sources.
Deep-learning approaches attempt to mitigate by instead learning the reconstruction mapping from data \citep{erkoc2026}, which may prove to be worth investigating in future revisions.
